\documentclass[11pt,a4paper]{article}

\usepackage[T1]{fontenc}
\usepackage[utf8]{inputenc}
\usepackage{lmodern}
\usepackage{amsmath}
\usepackage{amssymb}
\usepackage{amsthm}
\usepackage{booktabs}
\usepackage{array}
\usepackage{listings}
\usepackage{microtype}
\usepackage[margin=1in]{geometry}
\usepackage[hidelinks]{hyperref}

\lstset{
  basicstyle=\ttfamily\footnotesize,
  columns=fullflexible,
  keepspaces=true,
  showstringspaces=false,
  xleftmargin=1.5em,
  frame=none,
  breaklines=true
}

\theoremstyle{plain}
\newtheorem{theorem}{Theorem}
\newtheorem{lemma}[theorem]{Lemma}
\newtheorem{proposition}[theorem]{Proposition}
\newtheorem{corollary}[theorem]{Corollary}
\theoremstyle{definition}
\newtheorem{definition}[theorem]{Definition}
\newtheorem{hypo}{H}
\renewcommand{\thehypo}{\arabic{hypo}}
\theoremstyle{remark}
\newtheorem{remark}[theorem]{Remark}

\newcommand{\Reach}{\mathrm{Reach}}
\newcommand{\wid}{w}
\newcommand{\nmin}{n_{\min}}
\newcommand{\nstab}{n_{\mathrm{stab}}}
\newcommand{\cube}{\mathrm{cube}}
\newcommand{\Succ}{\mathrm{Succ}}
\newcommand{\Prune}{\mathrm{Prune}}
\newcommand{\canon}{\mathrm{canon}}

\title{A Machine-Checked Cutoff Theorem for the Census of an\\
       Exhaustive Combinatorial-Optimality Search}
\author{}
\date{}

\begin{document}
\maketitle

\begin{abstract}
\noindent
Exhaustive searches for combinatorial optima are re-run from scratch at every
instance size, and their cost is dominated by the set of states they store.
We identify a class of such searches --- \emph{graded ambient searches} --- for
which that stored set stabilises: for every level $\ell$ and every ambient
$n$ above a cutoff $\nmin(\ell)$, the census at $n$ is determined by one
computation at $\nmin(\ell)$ together with an explicit \emph{affine} growth
term of slope exactly one. The theorem, its level-agreement clause, and a
conditional polynomial refinement are machine-checked in Isabelle/HOL with two
non-vacuous interpretations. The closest prior art is dynamic cutoff detection
for parameterized concurrent programs; we state what survives the comparison.
We instantiate the class in three domains, use it to re-certify $S(11)=35$
against an unmodified verified checker, and report the two measured walls at
which the method's practical reach ends.
\end{abstract}

%======================================================================
\section{Introduction}
%======================================================================

An exhaustive search for a combinatorial optimum spends its budget on
\emph{storage}. Branch-and-bound with memoisation, dynamic programming over
canonical states, and the depth- or size-optimality searches for sorting
networks~\cite{Knuth73,BZ14,CCFS16,Harder20} all share the same shape: a root
state, a transition relation, a
bound oracle, and a memo table whose size is the binding constraint. The memo
table is a function of the instance parameter $n$, and the conventional
assumption is that each $n$ demands its own computation. When the search is
already at the edge of a machine at $n$, the computation at $n+1$ is simply
out of reach.

This paper identifies a structural condition under which that assumption is
false. Let a search carry, besides its instance parameter $n$ --- which we call
the \emph{ambient} --- an intrinsic \emph{width} $\wid(s) \le n$ on states,
and let its cost budget be reparameterised as a \emph{level} $\ell$ above a
forced prefix of the search whose cost $C(n)$ is known in closed form. Write
$\Reach(n,\ell)$ for the census: the set of states the search stores at ambient
$n$ and level $\ell$. Under five hypotheses on naming, transitions, grading and
geodesics, the census satisfies an exact identity
\[
\Reach(n,\ell)\;=\;\Reach(\nmin(\ell),\ell)\;\uplus\;
\{\,\mathrm{root}_w : \nmin(\ell) < w \le n\,\},
\]
so that its cardinality grows by exactly one per unit of ambient above the
cutoff. One computation determines the census at every larger instance size,
forever. We prove this in two logically independent halves --- an
\emph{ambient collapse} that needs no cutoff at all, and a \emph{chain
collapse} that supplies one --- and we machine-check both.

\paragraph{Prior art, conceded first.}
This is a \emph{cutoff theorem}, and cutoffs are a mature named concept in
parameterized verification~\cite{EN95,EK00,Nam07,AJK16,Jaber20,HS20,BN23,BloemEtAl15}.
The closest prior art is not a distant analogy but a direct one.
Kaiser, Kroening and Wahl~\cite{KKW10}, Definition~4, define a cutoff as
``a number $c \in \mathbb{N}$ such that, for all $n \ge c$, $R_n = R_c$'',
where $R_n$ is a \emph{set of reachable states} of an $n$-thread program ---
not the truth value of a property. An earlier version of the present work
claimed to be the first cutoff theorem about a set rather than a property.
That claim is false and we retract it here. Three distinctions survive, and we
claim only these. (i)~$R_n$ in~\cite{KKW10} is a set of \emph{abstract} thread
states under a fixed finite abstraction; our census is the concrete set of
states an exhaustive optimality search actually stores, at full fidelity.
(ii)~\cite{KKW10} concludes exact equality with no growth term; we give an
affine growth law with slope exactly one, and, conditionally, a polynomial
one. (iii)~Their cutoff is \emph{dynamically detected} by a sufficient runtime
test during the analysis; ours is an \emph{a-priori} consequence of the
hypotheses, computable before the search is launched. Two adjacent bodies of
work are not prior art: symmetry and partial-order reduction shrink an explored
state space by a quotient or a representative
selection~\cite{CEFJ96,ID96,Valmari91}, and work on models of dynamic
programming and dominance caching in branch-and-bound~\cite{BODI11,CGS24}
concerns expressiveness and pruning --- neither asserts that the object
explored at one parameter equals the object explored at another. We use
``cutoff'' in the
tight-cutoff sense of Emerson--Namjoshi~\cite{EN95,EN03} and
Au{\ss}erlechner--Jacobs--Khalimov~\cite{AJK16}, with the census of the explored
state set in place of the truth of a temporal property; not in the
Emerson--Kahlon sense~\cite{EK00}. Namjoshi's completeness result for
invariance~\cite{Nam07} is the other landmark of the a-priori side: there the
cutoff exists but finding it is undecidable, whereas ours is read off the
hypotheses.

\paragraph{Scope, conceded second.}
The class defines the census as a ball in a level-independent weighted
distance. The engine we instantiate it on stores a \emph{bound-gated subset} of
that ball, because its recursion reaches successor expansion only when a
Huffman-style bound test passes. The two disagree already at level~1. The
consequence is stated once and enforced throughout: the sorting-network
identity is proved directly from the gate, not as a corollary of the abstract
theorem, and the claim that the sorting-network search satisfies the five
hypotheses is \emph{not} machine-checked. It rests on an exhaustive audit of
the implementation, reported as such.

\paragraph{Origin.}
The class was extracted from a concrete question. A 1972 lower-bound argument
for sorting networks is the sole published support for a value of $S(13)$ that
circulates in the two most-consulted public tables, and a companion
paper~\cite{Paper1} shows that the argument does not establish it. Auditing the
computational route to a replacement meant asking what a size-optimality search
at $13$ channels would have to store, which meant relating it to the same
search at $11$ channels; that is where the collapse identity came from. The
audit does not recur here beyond Section~\ref{sec:instance}, where the
sorting-network search is one of three instances.

\paragraph{Road map.}
Section~\ref{sec:class} defines the class and states the five hypotheses,
including the one --- width convexity --- that was invisible in the concrete
statement because it happens to be free there.
Section~\ref{sec:theorem} states both halves of the theorem, the
level-agreement clause and the growth law, separating the affine law we prove
from the polynomial law we do not. Section~\ref{sec:mech} reports the Isabelle
development, the two interpretations that make it non-vacuous, three
corrections the mechanization forced on our own hypotheses, and what is not
mechanized. Sections~\ref{sec:instance} and~\ref{sec:domains} give three
instances; Section~\ref{sec:s11} the certified computation;
Section~\ref{sec:boundary} the measured boundary; Section~\ref{sec:trust} the
trusted base. Everything a reader must audit is in the paper; the theory file,
the engine, the certificate and the verifiers are in the
deposit~\cite{Artifact}.

%======================================================================
\section{Graded ambient searches}\label{sec:class}
%======================================================================

\begin{definition}[graded ambient search]\label{def:gas}
A \emph{graded ambient search} is a tuple $(S, \wid, \to, \gamma,
\{\mathrm{root}_n\})$ in which $S$ is a set of canonical states,
$\wid : S \to \mathbb{N}$ assigns each state an intrinsic \emph{width},
$\to_n\, \subseteq S \times S$ is a transition relation for each ambient $n$,
$\gamma$ assigns each edge a non-negative integer \emph{grade}, and
$\mathrm{root}_n \in S$ is the ambient-$n$ root. Crucially $S$ itself does not
depend on $n$. Writing $d_n$ for the $\gamma$-weighted distance in $\to_n$ and
$C(n)$ for the free-chain offset of~H3 below, the object of study is the
\emph{census}
\[
\Reach(n,\ell) \;=\; \{\, s \in S : d_n(\mathrm{root}_n, s) \le C(n) + \ell \,\},
\]
where $\ell$ is the \emph{level} and, for a width value $v$,
$W(v,\ell) = \{s \in \Reach(v,\ell) : \wid(s) = v\}$ is the width-$v$
\emph{stratum}.
\end{definition}

Three features are deliberate. The census is the set of states an exhaustive
search \emph{stores}, not the abstract solution set: it is what costs money.
The width is intrinsic to the state, not to the run that found it, so
canonicalisation must quotient away unused dimensions. And the budget is
reparameterised: a search whose first $C(n)$ units of progress are forced is
studied at \emph{level} $\ell = L - C(n)$, not at total cost $L$. Without that
shift the ambient dependence sits entirely in the offset and no identity is
visible.

\subsection{The hypotheses}

\begin{description}\itemsep2pt
\item[H0 (ambient-free naming).] The canonical name of a state is a function of
  the state alone; canonicalisation quotients away unused dimensions, so a
  state has one and the same name at every ambient $n \ge \wid(s)$.
\item[H1 (ambient freedom, stratified).] For every $s$ and all
  $\wid(s) \le m \le n$,
  $\{\,t : s \to_n t,\ \wid(t) \le m\,\} = \{\,t : s \to_m t\,\}$.
  Raising the ambient may only add transitions to \emph{wider} states.
\item[H2 (ambient-free grading).] $\gamma(s \to t)$ is a function of $(s,t)$
  alone.
\item[H3 (root tower).] $\mathrm{root}_n \to_n \mathrm{root}_{n-1}$ with grade
  $\delta(n)$, and $C(n) = \sum_{k \le n} \delta(k)$ is the \emph{free chain}.
  The \emph{degenerate} case $\mathrm{root}_n = \mathrm{root}_0$, $C \equiv 0$,
  is explicitly allowed.
\item[H4 (width convexity).] For all $m \le n$ and every $s$ with
  $\wid(s) \le m$, every ambient-$n$ walk $\mathrm{root}_n \rightsquigarrow s$
  of grade $g$ can be replaced, at grade no greater than $g$, by the root tower
  $\mathrm{root}_n \rightsquigarrow \mathrm{root}_m$ of grade $C(n)-C(m)$
  followed by a walk confined to states of width $\le m$.
\item[H5 (gated tail).] There is a computable $\nmin(\ell)$ such that for every
  $v > \nmin(\ell)$ the stratum $W(v,\ell)$ is the single chain state
  $\mathrm{root}_v$.
\end{description}

H0 and H2 are signature constraints rather than hypotheses, and in the
mechanization they become exactly that (Section~\ref{sec:mech}). H1 says the
ambient is a ceiling and nothing else: it may open routes into wider states,
but it may not alter the width-$\le m$ part of the transition structure. H3
names the forced prefix and may be trivial.

H4 is where the mathematics is. It says that no detour through high-width
states is ever a shortcut. Where width can only decrease along an edge, H4 is
free and therefore invisible: our own first statement of the theorem, phrased
for sorting networks, did not contain it, because there $\Succ$ preserves width
and $\Prune$ lowers it by one, so every geodesic descends monotonically. The
moment an edge can \emph{raise} the width, a large-ambient search acquires
routes unavailable at a smaller ambient that nevertheless end inside the
width-$\le m$ stratum, and H4 becomes the genuine claim that those routes never
help. That is why the second domain of Section~\ref{sec:domains} was chosen to
be one in which width is neither monotone up nor down: H4 is the hypothesis
that can fail, and it deserved an adversarial test.

The form of H4 printed above is a \emph{decomposition}: tower, then confinement.
It is not the form we first published, and the difference is not cosmetic; see
Section~\ref{sec:corrections}.

H5 is strong and rare, and is stated above at the small ambient $v$ itself
rather than as a property of the large search: weaker as a hypothesis, stronger
as a theorem. It forces $\wid(\mathrm{root}_v) = v$ above the threshold, so it
cannot hold with a degenerate root tower unless the strata are empty.

\subsection{Which half needs what}

The hypotheses split the theorem in two, and the split is the main structural
finding of the abstraction.

\begin{center}
\begin{tabular}{lcccccc}
\toprule
proof step & H0 & H1 & H2 & H3 & H4 & H5 \\
\midrule
geodesic confinable to width $\le m$          &     &     &     &     & yes &     \\
initial segment is the root tower, grade $C(n)-C(m)$ &     &     &     & yes &     &     \\
remaining edges exist at ambient $m$          &     & yes &     &     &     &     \\
grades agree                                  &     &     & yes &     &     &     \\
states carry the same names at both ambients  & yes &     &     &     &     &     \\
converse direction                            &     & yes & yes & yes &     &     \\
the tail above $\nmin(\ell)$ is a chain       &     &     &     & yes &     & yes \\
\bottomrule
\end{tabular}
\end{center}

Part~I (ambient collapse) needs H0--H4; Part~II (chain collapse) additionally
needs H5 \emph{and} needs H3 to be non-degenerate. The useful statement is not
that Part~I is independent of H3 --- it is not --- but that Part~I holds with a
\emph{degenerate} root tower, in which case H3 is vacuous. The split is
therefore between the hypotheses that make the tail trivial
(H3-non-degenerate together with H5) and everything else.

%======================================================================
\section{The theorem and the growth law}\label{sec:theorem}
%======================================================================

\begin{theorem}[Ambient Collapse]\label{thm:partI}
Assume H0--H4. Then for all $m \le n$ and every level $\ell$,
\[
\Reach(n,\ell) \cap \{\, s : \wid(s) \le m \,\} \;=\; \Reach(m,\ell),
\]
and $\mathrm{level}_n(s) = \mathrm{level}_m(s)$ for every $s$ of width
$\le m$, where $\mathrm{level}_n(s)$ is the least $\ell$ with
$s \in \Reach(n,\ell)$.
\end{theorem}

\begin{proof}[Proof sketch]
Fix $s$ with $\wid(s) \le m \le n$ and an ambient-$n$ walk from
$\mathrm{root}_n$ to $s$ of grade $g \le C(n) + \ell$. By H4 replace it, at no
greater grade, by the root tower down to $\mathrm{root}_m$ --- of grade exactly
$C(n) - C(m)$ by H3 --- followed by a walk confined to width $\le m$ of grade
$g'$ with $(C(n)-C(m)) + g' \le g$. Every edge of the confined part joins states
of width $\le m$, so by H1 it is an edge of $\to_m$, by H2 carries the same
grade, and by H0 its endpoints have the same names at both ambients; hence
$g' \le C(m) + \ell$. Conversely, every edge of a walk at ambient $m$ joins
states of width $\le m$, so by H1 read in the other direction it is an edge of
$\to_n$, and prepending the root tower gives an ambient-$n$ walk of grade
$C(n)-C(m)+g'$. Equality of levels follows, hence of the width-$\le m$ parts.
\end{proof}

\begin{corollary}[stratification]\label{cor:strat}
Under H0--H4, $\Reach(n,\ell) = \biguplus_{v \le n} W(v,\ell)$, where each
stratum is a function of $(v,\ell)$ alone and carries no dependence on $n$.
\end{corollary}

This is the form a referee should be shown. Theorem~\ref{thm:partI} is an
identity between two sets that both mention $n$; Corollary~\ref{cor:strat} says
the ambient has been eliminated --- the census is a running sum of $n$-free
pieces --- and it follows in one line by applying Theorem~\ref{thm:partI} at
$m := \wid(s)$.

\begin{theorem}[Chain Collapse]\label{thm:partII}
Assume H0--H5, with H3 non-degenerate. Then for all $n \ge \nmin(\ell)$,
\[
\Reach(n,\ell) \;=\; \Reach(\nmin(\ell),\ell) \;\uplus\;
\{\, \mathrm{root}_v : \nmin(\ell) < v \le n \,\},
\]
and if $\Reach(\nmin(\ell),\ell)$ is finite then
$|\Reach(n,\ell)| = |\Reach(\nmin(\ell),\ell)| + (n - \nmin(\ell))$.
\end{theorem}

Theorem~\ref{thm:partII} is the cutoff proper: the census at every ambient
follows from one computation at $\nmin(\ell)$, and the correction term is an
explicit closed form rather than an error bound.

\subsection{Three tail regimes}

H5 is the strongest hypothesis in the list, and the interesting question is
what happens without it. Part~I alone forces the census to be a running sum of
$n$-free strata, so $|\Reach(n,\ell)|$ for large $n$ is determined entirely by
the shape of $W(\cdot,\ell)$ above the population front. Across the three
domains below, exactly three regimes occur.

\begin{center}
\begin{tabular}{llll}
\toprule
regime & strata above the front & $|\Reach(n,\ell)|$, large $n$ & instance \\
\midrule
T-empty   & nothing                 & constant             & Boolean chains \\
T-trivial (= H5) & one state per width & $\mathrm{const} + n$ & sorting networks \\
T-poly    & a polynomial in $w$     & polynomial in $n$    & prefix reversal \\
\bottomrule
\end{tabular}
\end{center}

\begin{corollary}[census growth law]\label{cor:growth}
Assume H0--H4, that every stratum is finite, and that $W(\cdot,\ell)$ agrees
with a polynomial of degree $\le d$ for all $v \ge v_0$. Then
$|\Reach(n,\ell)|$ agrees, for $n \ge v_0$, with a polynomial in $n$ of degree
$\le d+1$; consequently one run at ambient $v_0 + d + 1$ determines the census
at every larger ambient exactly.
\end{corollary}

Corollary~\ref{cor:growth} is what makes the method survive the loss of H5, and
in practice T-poly is the regime it most often lands in.

\subsection{Affine, not polynomial: a distinction we must not blur}

Three different statements travel under the phrase ``polynomial growth law''
and they have three different statuses. We separate them once and keep them
separate.

\emph{Growth in the ambient, gated case (PROVEN).} Under H0--H5 the census is
eventually affine in $n$ with slope exactly $1$: Theorem~\ref{thm:partII},
machine-checked. Degree one is a polynomial, so ``polynomial growth law'' is not
false, but it oversells and obscures the sharp part; what distinguishes this
from~\cite{KKW10} --- exact equality with no growth term at all --- is fully
carried by ``affine with slope exactly one''. For the sorting-network search we
have checked it against every census figure in the artifact: at levels
$1,2,3,4$ the census over ambients $6$--$13$, $7$--$13$, $9$--$13$ and
$9$--$13$ is an arithmetic progression of common difference exactly $1$, and at
levels $1$--$5$ the width census at ambient $13$ holds exactly one state at
every width strictly above $\nmin(\ell)$.

\emph{Growth in the ambient, polynomial case (CONDITIONAL; machine-checked as
an implication).} Corollary~\ref{cor:growth} is a mechanized theorem, but its
hypothesis --- that the strata are eventually polynomial --- is a measured fact
in the one domain where it applies, not a derived one, and there we prove only
the matching upper bound (Section~\ref{sec:domains}). ``A polynomial growth law
for the census'' is defensible only as this conditional, and only for that
domain.

\emph{Growth in the level (NOT A LAW).} For the sorting-network search
$|\Reach(\nmin(\ell),\ell)|$ at $\ell = 1,\dots,6$ is
\[
33,\quad 445,\quad 24{,}199,\quad 207{,}995,\quad 50{,}602{,}869,\quad
95{,}221{,}143,
\]
with successive ratios $13.48,\ 54.38,\ 8.60,\ 243.29,\ 1.88$ --- a spread of
$241\times$. No forward difference of orders one through five vanishes anywhere
in the measured range. The alternation is structural: the ratio is large
exactly when $\nmin$ jumps and small when it does not. Nothing in this paper
bounds growth in the level, and any sentence about polynomial growth of the
census must say \emph{in the ambient}.

%======================================================================
\section{The mechanization}\label{sec:mech}
%======================================================================

The abstract theorem is machine-checked in Isabelle/HOL~\cite{Isabelle}. The
development is one file, \texttt{Ambient\_Collapse.thy}, $1{,}095$ lines, built
under Isabelle2020 in its own session in about two seconds on top of a prebuilt
HOL heap. It contains no \texttt{sorry}, no \texttt{oops}, no
\texttt{axiomatization} and no \texttt{quick\_and\_dirty}; the build script
greps for all four and refuses to invoke Isabelle if any is present. The file
\texttt{imports Main} and nothing else --- in particular it never touches the
verified checker theories of Section~\ref{sec:s11}, so the two formal artifacts
are independent.

\subsection{The locale}

Walks are formalised by an inductive predicate \texttt{pathg E G s t g} --- a
walk from $s$ to $t$ of grade exactly $g$ --- and a variant \texttt{pathgB}
carrying a predicate every intermediate state must satisfy, used for width
confinement. The census is defined by an \emph{inequality} on the grade,
\begin{lstlisting}
Reach n l = {s. EX g. pathg (edge n) grade (root n) s g & g <= C n + l}
\end{lstlisting}
which is definitionally the ball $\{s : d_n(\mathrm{root}_n,s) \le C(n)+\ell\}$
whenever the distance exists and is total when it does not, so the weighted
distance never has to be constructed and no partiality enters.

H0 and H2 are discharged \emph{structurally} and cannot be violated by any
interpretation: H0 is the fact that states inhabit a single type \texttt{'s}
not mentioning $n$, so a state has one name at every ambient by construction,
and H2 is the fact that \texttt{grade} has type
\texttt{'s $\Rightarrow$ 's $\Rightarrow$ nat}. These are signature constraints
rather than hypotheses, and they are the honest formal content of
``canonicalisation quotients away the unused dimensions''. The rest are locale
assumptions:

\begin{lstlisting}
locale graded_ambient_search =
  fixes width :: "'s => nat"  and edge :: "nat => 's => 's => bool"
    and grade :: "'s => 's => nat"  and root :: "nat => 's"
    and C :: "nat => nat"
  assumes H1: "[| width s <= m; m <= n |] ==>
                 {t. edge n s t & width t <= m} = {t. edge m s t}"
    and H3_width: "width (root m) <= m"
    and H3_mono:  "m <= n ==> C m <= C n"
    and H3_tower: "m <= n ==> pathg (edge n) grade (root n) (root m) (C n - C m)"
    and H4: "[| m <= n; width s <= m; pathg (edge n) grade (root n) s g |] ==>
               EX g'. pathgB (edge n) grade (%t. width t <= m) (root m) s g'
                      & (C n - C m) + g' <= g"

locale gated_ambient_search = graded_ambient_search +
  fixes nmin :: "nat => nat"
  assumes H5: "nmin l < v ==> stratum v l = {root v}"
\end{lstlisting}

H1 is transcribed verbatim as a set equality; the mechanization extracts its
three consequences separately, one of which --- that an ambient-$m$ edge out of
a width-$\le m$ state lands in width $\le m$ --- is implied by the equality but
was never stated in our prose, and is used twice.

\subsection{What is proved}

The following are theorems of the locale, stated here exactly as they appear in
the file.

\begin{lstlisting}
theorem ambient_collapse:    "m <= n ==> Reach n l Int {s. width s <= m} = Reach m l"
corollary level_agree:       "[| m <= n; width s <= m |] ==> level n s = level m s"
theorem census_stratified:   "Reach n l = (UN v<=n. stratum v l)"
theorem chain_collapse:      "nmin l <= n ==>
                                Reach n l = Reach (nmin l) l
                                            Un root ` {v. nmin l < v & v <= n}"
theorem chain_collapse_card: "[| nmin l <= n; finite (Reach (nmin l) l) |] ==>
                                card (Reach n l)
                                  = card (Reach (nmin l) l) + (n - nmin l)"
theorem census_affine:       "[| !!v. finite (stratum v l); nmin l <= n |] ==>
                                fdiff (%n. census n l) n = 1"
(* stratum v l = {s : Reach v l. width s = v};  fdiff f n = f (Suc n) - f n *)
\end{lstlisting}

Corollary~\ref{cor:growth} is mechanized as \texttt{census\_growth\_law} in
vanishing-forward-difference form, then converted by a supporting lemma
\texttt{newton\_expansion} --- vanishing $(d{+}1)$-st differences imply
agreement with $\sum_{j \le d+1} c_j \binom{t}{j}$ --- proved constructively by
telescoping and the hockey-stick identity. Two reasons for the difference form.
First, vanishing forward differences are what the diagnostic tooling measures,
so the mechanized statement and the measured statement coincide. Second, it
makes \texttt{census\_affine} literally the $d=0$ case: T-trivial and T-poly
are one theorem at two degrees.

\subsection{Non-vacuity: two interpretations}

Locale theorems are worthless if the assumptions are contradictory, so both
locales are interpreted. Neither interpretation formalizes a real domain; both
are the smallest honest models that exercise the hypotheses, and both discharge
every obligation.

\emph{\texttt{Words}} --- states are \texttt{nat list}, the ambient is the
alphabet size, the width of a word is the least alphabet containing it, and the
two moves, each of grade one, are appending a letter $a < n$ and dropping the
last letter. The tower is degenerate. The mechanization checks that
$\texttt{Reach } 2\ 1 \subsetneq \texttt{Reach } 3\ 1$ --- so the ambient
really matters --- and that there are edges which \emph{raise} the width and
edges which \emph{lower} it. H4 therefore has genuine content here and is
proved, not assumed away: because every move changes the length by exactly one
and costs exactly one, a word of length $k$ cannot be reached in fewer than $k$
moves however far the walk climbs in between. That is the ``no high-width
detour is a shortcut'' statement, discharged.

\emph{\texttt{Tower}/\texttt{Gated}} --- a two-constructor datatype whose chain
$\mathrm{Tow}(n) \to \mathrm{Tow}(n-1) \to \cdots$ has grades $\delta(k)$,
giving a non-degenerate free chain, and in which each $\mathrm{Tow}(k)$ has
exactly one off-chain successor of the \emph{same} width, reachable only at
grade $C(k) + D(k) + 1$. With $D$ monotone and unbounded, H5 holds; the
cardinality identity and $\texttt{fdiff}(\lambda n.\ \texttt{census } n\ l) = 1$
above the cutoff are both checked. This is the abstract skeleton of the
sorting-network instance and the honest bridge for the scope gap of
Section~\ref{sec:instance}: it reproduces the bound gate by \emph{pricing} the
off-chain edge, turning a level-dependent control-flow test into a
level-independent grade. It is not a formalization of that search: here the
gate is postulated, and in the real engine it is derived.

The \texttt{Words} interpretation also earns its keep numerically. Its
width-$v$ stratum at level $\ell$ is $\sum_{k \le \ell}(v^k - (v-1)^k)$, of
degree $\ell-1$ with leading coefficient exactly $\ell$, and its sphere at
distance $\ell$ has size $n^\ell$. Those are precisely the degree and leading
coefficient the prefix-reversal domain exhibits
(Section~\ref{sec:domains}) --- produced by a model with no prefix-reversal
content whatsoever. We draw the deflationary conclusion in place rather than
leave it to a referee.

\subsection{Three corrections the mechanization forced}\label{sec:corrections}

Mechanizing a theorem one believes is a way of finding out what one actually
believes. Three of our own published hypotheses did not survive.

\emph{(1) Part~I needs H3.} We had asserted that Part~I uses neither H3 nor H5.
That is false, and it contradicted our own step table. Part~I relates the
ambient-$n$ distance from $\mathrm{root}_n$ to the ambient-$m$ distance from
$\mathrm{root}_m$; without a hypothesis relating the two roots and pricing the
passage between them, the distances are incomparable. In the mechanization H3
appears twice: as \texttt{H3\_tower} in the converse direction, and inside the
statement of H4 in the forward direction, where the $C(n)-C(m)$ term is what
accounts for the tower prefix. What survives is the statement at the end of
Section~\ref{sec:class}: Part~I is independent of H5 and holds with a
\emph{degenerate} root tower.

\emph{(2) H4 as originally printed is ill-formed.} Our published H4 read: for
every $s$ with $\wid(s) \le m$ and every $n \ge m$, some $\gamma$-geodesic
$\mathrm{root}_n \rightsquigarrow s$ passes only through states of width
$\le \max(m, \wid(\mathrm{root}_m))$. When the tower is non-degenerate,
$\wid(\mathrm{root}_n) = n > m$ and $\mathrm{root}_n$ lies on every geodesic
out of itself, so the condition is unsatisfiable in exactly the case in which
it is interesting --- which is the sorting-network case. It is satisfiable in
the two degenerate-tower domains, which is why our Python hypothesis-checkers
never caught it. The repair is the decomposition form of
Section~\ref{sec:class}, fusing the first two steps of the informal proof into
one hypothesis; the $\max$ form is not repairable by adjusting the bound.

Being candid about what this costs: H4-as-decomposition now carries the
mathematical weight of Part~I. The mechanization proves the \emph{transport}
(a width-confined ambient-$n$ walk is an ambient-$m$ walk of the same grade,
and conversely) and the \emph{bookkeeping} (the levels line up under
$C(n)-C(m)$). It does not prove H4 for any real domain. In the two
interpretations H4 is discharged, non-trivially in \texttt{Words}; in sorting
networks it is asserted to follow from width closure, and in prefix reversal it
is verified over a finite range, not proved.

\emph{(3) H5 belongs at the small ambient.} We had stated H5 as a property of
the stratum inside the large search. Stating it at ambient $v$ itself and
letting Part~I lift it is weaker as a hypothesis and stronger as a theorem, and
it makes the mechanized Part~II a derivation rather than a restatement. It also
lets us \emph{derive} what we had only observed: the only domain with a
non-degenerate root tower is the only domain with a trivial tail.

\subsection{What is not mechanized}\label{sec:notmech}

Stated flatly, because the distinction is the entire value of the section.

\begin{enumerate}\itemsep2pt
\item \textbf{That the sorting-network search is an instance is not
  mechanized.} It rests on an exhaustive reading of the engine's recursion,
  plus a second reviewer's audit (Section~\ref{sec:instance}). What changed
  with the mechanization is the \emph{shape} of the residual risk: before, an
  abstract theorem proved on paper resting on an unverified implementation
  property; now, a machine-checked theorem plus an unverified claim that one
  concrete search satisfies five stated properties. The second obligation is
  smaller and far better localised. We do not say the collapse theorem ``is
  verified for sorting networks''.
\item \textbf{The class definition does not literally model that engine.} The
  census of Definition~\ref{def:gas} is a ball; the engine stores a bound-gated
  subset of one. Section~\ref{sec:instance} gives the level-1 disagreement
  explicitly.
\item \textbf{The two further domains are not formalized.} They are checked by
  a Python harness, with a brute-force validator for one of them; formalizing
  the prefix-reversal domain would require \emph{proving} H4 for it, which is
  open.
\item \textbf{Finiteness is an explicit hypothesis, not a consequence.} The
  three cardinality theorems all assume the strata finite; nothing in H0--H5
  implies it, and it is discharged in both interpretations.
\end{enumerate}

%======================================================================
\section{The sorting-network instance}\label{sec:instance}
%======================================================================

The concrete search is Harder's size-optimality dynamic program for sorting
networks~\cite{Harder20}, the engine that settled $S(11)=35$ and $S(12)=39$;
it is the size-optimality counterpart of the SAT- and
symmetry-break-driven searches that settled the smaller
cases~\cite{CCFS16,CCS15,BZ14,EM14}. A
state is a canonical output set of width $k$; the two moves are applying a
comparator ($\Succ$, width-preserving) and conditioning on an extremal channel
and deleting it ($\Prune$, width-lowering by one); the bound oracle combines a
Huffman-style rule~\cite{Huffman52} over the pruned children with successor
expansion. The ambient is the channel count.

The semantic half of the instantiation is \emph{not} a mathematical novelty,
and we say so first. Harder defines the objective for a sequence set with no
reference to any ambient, so ``a bound on $X$ is valid at every ambient
$n \ge \wid(X)$'' is immediate from his definition. What is ours is an
\emph{implementation audit}: the claim that the engine actually respects that,
established by reading every loop bound, table index and ordering key inside
the recursion and confirming each is the current state's own channel count or a
stored bound; a second reviewer given the same code as a refutation target; and
the identification of the single place where it fails.

That place is the on-line subsumption index, whose default width policy is
keyed to the \emph{root} width, at widths $\{n-3, n-2\}$ --- different absolute
widths at every ambient --- and which is read inside the recursion. With it
enabled under that policy, H1 is false. The repair is one line of
configuration, not of code: pin the policy to absolute widths. The collapse
check is what found it, and an identical synthetic leak in a control domain
(Section~\ref{sec:domains}) breaks the collapse immediately and visibly.

Given the audit, the instance rests on three facts about the root, each
machine-checked against dumps of the engine's stored states. Every channel of
the full cube is extremal for both polarities, so $\Prune(\cube_k)$ is the
single state $\cube_{k-1}$ and the Huffman rule at $\cube_k$ combines $k$
identical children: the free chain is the chain of full cubes and
$C(k) = C(k-1) + \lceil \log_2 k\rceil$. The symmetric group acts transitively
on unordered pairs and fixes $\cube_k$ setwise, so $\Succ(\cube_k)$ is a single
state. The Huffman ceiling at $\cube_k$ is $S(k-1) + \lceil \log_2 k \rceil =
C(k) + D(k-1)$ where $D(n) = S(n) - C(n)$, and successor expansion at the root
is not reached below that value; hence H5 holds with
$\nmin(\ell) = \min\{\, n : D(n) \ge \ell \,\}$ and the tail is one state per
width. Note that $C(n) = \sum_{k\le n}\lceil\log_2 k\rceil$ is Knuth's
binary-insertion number~\cite{Knuth73}: classical, and what is ours is only its
identification as the cost of a forced prefix and the reparameterisation
$\ell = L - C(n)$.

\paragraph{The scope gap, stated plainly.}
The engine's census is not the ball of Definition~\ref{def:gas}. Grade is one
per comparator, so the unique width-13 successor of $\cube_{13}$ lies at
distance $1$ and belongs to the ball at level~$1$; but the measured level-1
census at ambient $13$ contains exactly one state of width $13$, namely
$\cube_{13}$ itself. Ball and census disagree at level~1. The reason is the
gate: the recursion reaches successor expansion only when the Huffman test
passes, and that test depends on the level, which
Definition~\ref{def:gas} forbids. Consequently \emph{the abstract theorem is a
guide here, not a proof}. The sorting-network identity is proved directly from
the gate, and must not be presented as a corollary of
Theorem~\ref{thm:partII}.

\paragraph{The prediction, and the folklore objection.}
From the machine-checked table $D(3\ldots12) = 0,0,1,1,2,2,4,4,6,6$ we get
$\nmin(\ell) = 5,7,9,9,11,11,13$ for $\ell = 1,\dots,7$. This was derived
before the censuses were computed and then checked against them: the population
front of the ambient-13 level-$\ell$ census sits at width $5$, $7$, $9$ for
$\ell = 1,2,3$, and every width above the front holds exactly one state, at
lower bound $C(w)+\ell$. At level 5 --- $50{,}602{,}871$ states --- the front is
at width $11$, and widths $12$ and $13$ hold one state each.

The bounds these censuses bear on are also the ones circulating in the public
tables of sorting-network sizes~\cite{Dobbelaere25}, and lower bounds for
restricted classes remain an active target~\cite{Tozawa26,VV72,FK73}. The
deflationary explanation for any census stabilising in $n$ is the folklore
support bound: an object built from $L$ operations touches at most $2L$
coordinates, so the census cannot change once $n \ge 2L$. Here $L = C(n)+\ell$
with $C(n) = \Theta(n \log n)$, so the support bound demands $n \ge 2(C(n)+\ell)$,
which is false for every $n \ge 3$ and becomes more false as $n$ grows.

\begin{center}
\begin{tabular}{ccccc}
\toprule
level $\ell$ & $\nmin(\ell)$ & comparators $L = C(\nmin)+\ell$ & support bound $2L$ & $2L \le \nmin$? \\
\midrule
1 & 5  & 9  & 18 & no ($3.6\times$ over) \\
2 & 7  & 16 & 32 & no ($4.6\times$ over) \\
3 & 9  & 24 & 48 & no ($5.3\times$ over) \\
4 & 9  & 25 & 50 & no ($5.6\times$ over) \\
5 & 11 & 34 & 68 & no ($6.2\times$ over) \\
6 & 11 & 35 & 70 & no ($6.4\times$ over) \\
\bottomrule
\end{tabular}
\end{center}

The support bound never certifies stabilisation at any ambient the search can
reach, at any level, and the gap widens with $n$. The cutoff here comes from
the gated root tower, which is a different mechanism.

%======================================================================
\section{Two further domains}\label{sec:domains}
%======================================================================

A method claim needs instances that were not the source of the abstraction. We
ran two under a pre-registration protocol: compute the small instance, derive
the prediction \emph{from that file alone}, hash it, record the hash, and only
then compute the larger ambients. The check is on the SHA-256 of the sorted key
list, not on counts, so a collapse holding numerically but not state by state
would fail.

\paragraph{Boolean chains --- a positive control.}
A straight-line program over $\{\wedge,\vee,\oplus\}$ starts from $n$
projections and adds gates; the state is the multiset of gate outputs, the
level is the gate count, canonicalisation drops variables outside the essential
support of every computed function, and the width is the number of essential
variables. Width is \emph{non-decreasing} along an edge --- the opposite of the
sorting-network case, so Part~I's proof runs in the other direction --- and the
root has width $0$, so H3 is degenerate. Each gate introduces at most two fresh
variables, giving $\nmin(\ell) = 2\ell$ with an \emph{empty} tail. Every
pre-registered prediction held key for key: at level~3, ambients $2$--$10$, all
$36$ pairs $m<n$ pass with symmetric difference zero and no level
disagreements, the census being a step function saturating at $541$ from
ambient $6$; at level~4 the same at $11{,}263$ states from ambient $8$. A raw
enumeration with no canonicalisation and no modelling shortcut reproduces the
search exactly and shows the canonical form to be a \emph{complete} invariant,
so the collapse is not an artefact of a too-coarse quotient.

We grade this domain honestly: $\nmin(\ell) = 2\ell$ \emph{is} the folklore
support bound. It demonstrates that the framework is correct and that Part~I
holds with a degenerate tower and an upward-monotone width --- a positive
control, not independent evidence that census cutoffs ever beat folklore.

\paragraph{Prefix reversal --- the real test.}
A pancake flip $f_k$ reverses the first $k$ entries of a permutation of
$[n]$~\cite{Gates95,Bulteau15};
the level is the flip count, canonicalisation strips the maximal sorted suffix,
and the width is what remains. Two properties make this the adversarial case.
No support bound exists at all: a single flip $f_n$ produces a state of width
exactly $n$. And width is neither monotone up nor monotone down --- a flip past
the current width drags a sorted element to the front, and a later flip can
bring it back --- so a search at ambient $n$ genuinely has routes unavailable
at ambient $m < n$ that end inside the width-$\le m$ stratum. Part~I is
precisely the claim that no such detour is ever a shortcut, and that was the
pre-registered falsification target.

It was not falsified. At level~5, for all $45$ ambient pairs $m<n$ in
$3\ldots12$, the width-restricted censuses agree with symmetric difference zero
and no level disagreements on any shared key; a separate state-by-state check
of H1 against ambient $11$ found $0$ mismatches over $120$, $2{,}086$ and
$13{,}400$ states. At level~6, ambients $8,10,12,14$ --- up to $2{,}087{,}574$
states --- all $6$ pairs pass. With the two Boolean levels that is $93$ ambient
pairs $(m,n)$ with $m<n$ in total and no discrepancy anywhere. \textbf{H4 is verified, not
proved}: whether a high-width detour can be a shortcut at some larger level is
an open proof obligation, and the plausible argument (a flip past the width
displaces sorted elements that must all be restored) is not a proof.

Here H5 fails: the tail is $1, 4, 18, 96, 465, 1501, 3687, \dots$, and it is
exactly polynomial. The pre-registered fit --- degree $\ell-1$, leading
coefficient exactly $\ell$, valid for $w \ge \ell+1$, taken from widths
$\le 12$ only --- reproduced widths $13,14,15$ exactly at levels $3,4,5$, and
extrapolated to a level~6 that had not been computed when the prediction was
frozen: predicted $W(13,6) = 497{,}108$ from a degree-5 fit at $w=7\ldots12$,
observed $497{,}108$; sixth forward difference zero; leading coefficient $6$.
Summing the strata gives a statement about a classical object with no reference
to our machinery: \emph{the number of permutations of $[n]$ at prefix-reversal
distance exactly $\ell$ from the identity is a polynomial in $n$ of degree
exactly $\ell$ with leading coefficient exactly $1$}. We prove the upper bound
$(n-1)^\ell$; the matching lower bound is open. Eventual polynomiality of
pancake sphere sizes is in any case already known in the permutation-patterns
literature, which is the best outcome available for a method claim: the
abstract theorem predicted, from hypotheses checked mechanically in seconds, a
result proved independently elsewhere. It is a prediction of a known result,
not a new one --- and the deflationary reading of Section~\ref{sec:mech}
applies, since counting $\ell$-letter words over an $n$-letter alphabet already
gives that degree and that leading coefficient.

Two things did fail, and both are recorded. Our threshold formula
$\nstab(\ell) = \ell+1$, sharp for $\ell \le 6$, does not extend to $\ell = 7$,
where the onset is $\ell+2$; what survives unqualified is that the strata are
eventually polynomial of degree $\ell-1$ and that \emph{some} computable
$\nstab$ exists, which is all Corollary~\ref{cor:growth} needs. And a control
domain carrying a synthetic pruning heuristic keyed to the root width --- the
same shape as the sorting-network leak --- breaks the collapse at once, with
$285$ keys present at one ambient and absent at the next. The mathematics of
Part~I can be true while an implementation detail destroys its applicability;
any deployment of the method must run the check rather than assume it.

Read across the two domains: Part~I generalizes and Part~II does not.
\emph{``Chain collapse generalizes'' is false. ``Ambient collapse generalizes,
and chain collapse is the special case in which a gated root tower makes the
tail trivial'' is true}, and it is the stronger statement, because it says
\emph{why} the sorting-network search gets the clean identity. Neither new
domain has a non-degenerate root tower, and neither has a gate.

%======================================================================
\section{A certified census: \texorpdfstring{$S(11)=35$}{S(11)=35}}\label{sec:s11}
%======================================================================

The theorem's payoff is that a census computed once at $\nmin(\ell)$ stands for
every larger ambient. To exercise that we rebuilt the search from the published
description, re-derived $S(11) \ge 35$ from scratch, and had the certificate
checked by an unmodified extraction of the verified checker accompanying the
original work~\cite{Harder20}, in the lineage of the extracted certified
checkers of Cruz-Filipe and co-authors~\cite{CS15,CLS17} and resting on
Isabelle's code generator~\cite{HN10}. The point is not the value, which is
known: it is that an independently produced certificate is accepted by a
checker nobody here wrote, and that the census the run produces is the one the
theorem says determines every larger channel count.

The run was in four stages. Search, at ambient $11$ with the subsumption index
pinned to absolute widths $\{8,9\}$, took $71$~h $06$~m at a peak of
$13.08$~GB on a six-core Ryzen~3600 with 47~GB of memory, and returned the
bound $35$ with $95{,}221{,}143$ stored states --- reconstructed from the run
log as $103{,}343{,}397$ insertions less $8{,}122{,}254$ evictions. Prune-all
took $9$~h $45$~m at about $8.4$~GB. Certificate assembly took about $9$~h and
was killed twice by the out-of-memory killer, at an anonymous resident set of
$37.9$~GB on a machine with no swap; it completed on the third attempt only
after a 96~GB swap file was armed underneath the live process, so its true peak
is unknown and is by construction above $37.9$~GB. It emitted a certificate of
$2{,}442{,}317{,}348$ bytes ($2{,}329$~MiB), SHA-256
\texttt{672c433f}\ldots, containing $10{,}275{,}769$ steps. Verification, on an
Apple~M4, took $89$~m $24$~s at a peak of $4.35$~GB and returned
\texttt{Just (11,35)} with exit status~$0$. From the machine-checked
correctness proof of the certificate-checking routine we therefore get
$S(11) \ge 35$, matching the known upper bound.

\begin{center}
\small
\begin{tabular}{lllcc}
\toprule
stage & original~\cite{Harder20} & this work & memory & wall \\
\midrule
search      & 178~GiB (191~GB), 4~h 51~m & 13.08~GB, 71~h 06~m & $14.6\times$ less & $14.7\times$ more \\
prune-all   & 16~GiB, $\approx$2~d 5~h   & ${\sim}8.4$~GB, 9~h 45~m & $2.1\times$ less & $5.4\times$ less \\
gen-proof   & 54~GiB (58.0~GB), 19~h 02~m & $>37.9$~GB, ${\sim}9$~h & $1.5\times$ less & $2.1\times$ less \\
verify      & 6~GiB, 34~m                & 4.35~GB, 89~m       & $1.5\times$ less & $2.6\times$ more \\
certificate & 2{,}926~MiB                & 2{,}329~MiB         & $1.26\times$ smaller & --- \\
\midrule
total wall  & 77.5~h                     & 91.3~h              & ---               & $1.18\times$ more \\
\bottomrule
\end{tabular}
\end{center}

Three things about that table have to be said rather than inferred. First, the
$14.6\times$ is a \emph{search-stage} figure and nothing more; an earlier
version circulated as $13.6\times$, obtained by dividing $178$~GiB by
$13.08$~GB as though the units matched, and the corrected ratio is
$191.13/13.08 = 14.61$. Second, the memory was bought with time on weaker
hardware: the same stage is $14.7\times$ slower and end to end the pipeline is
$1.18\times$ slower --- a memory-for-time trade, not an improvement. Third, the
binding memory peak \emph{moved stages}: in the original it was search, here it
is certificate assembly, and the end-to-end improvement is only about
$1.5\times$ --- optimistic even so, since its denominator is a lower bound from
a failed attempt. Certificate generation, not search, is the hidden peak of
this pipeline and should be budgeted separately.

We also record what the run does \emph{not} establish, since the census is the
object the rest of the paper depends on. Only verification left a surviving
machine-readable artifact; the other three stages are recorded in prose and a
run log, and no timing output survives that would settle whether $13.08$ is
gigabytes or gibibytes. The wall-time column compares three machines and two
configurations and is not a controlled measurement. And the census figure that
Theorem~\ref{thm:partII} exports --- $|\Reach(13,6)| = 95{,}221{,}145$ --- is a
\emph{prediction} from the single ambient-11 measurement, not a measurement.

%======================================================================
\section{Where the method stops}\label{sec:boundary}
%======================================================================

A cutoff theorem tells you which computation to run. It does not tell you that
you can afford it, and in this programme the two answers diverge at the next
level. We report the divergence because it delimits the method.

Theorem~\ref{thm:partII} makes the level-6 census at ambient $13$ a corollary
of the ambient-11 run of Section~\ref{sec:s11}: $\nmin(6) = 11$, so no
computation at $13$ is needed. At level~7 the leverage vanishes, because
$\nmin(7) = 13$. There is no smaller ambient to collapse to; the theorem is
silent, and the cost is whatever the search costs. That is the method's
boundary in the sharpest form available. A cutoff at $\nmin(\ell)$ removes the
ambient as a cost driver and leaves the level as the only one --- and
Section~\ref{sec:theorem} showed that census growth in the level is governed by
no law we have.

\paragraph{What the cost actually tracks.}
The measurement that matters here is a diagnosis, not a projection, and it
corrects an assumption we had been making: extrapolating from state counts
assumes cost is proportional to stored states, and it is not. Between levels~5
and~6 of the certified run, stored entries grew $21.8\times$ and the
subsumption index population $20.4\times$, but index candidates examined grew
$297\times$ and the iteration wall $418\times$; index lookup accounted for
$79.4\%$ of user time. Two independent measurements --- the level-to-level
ladder, and a $25{,}613$-sample trace of the instantaneous state rate inside
level~6, which fell $6.32\times$ while the live set grew $10.3\times$ and was
still falling in the final hour --- agree that wall grows as $N^{1.79}$ to
$N^{1.94}$. So: \textbf{memory scales with the census; wall clock scales with
the index.} The census is the quantity our theorem controls; the wall is not,
and it is the faster-growing one. A
four-arm sweep of index configurations found a strict Pareto improvement over
both previously used settings --- $1.49\times$ faster and $4.2\times$ less
memory than an unindexed run, $2.02\times$ faster than the certified setting at
$1.86\times$ its memory --- which moves the operating point \emph{along} that
frontier rather than off it. Projecting each arm on its own measured ladder
puts a level-7 search at roughly $45$--$190$ days and $1.5$~TB resident against
a months-scale horizon and $143$~GB of memory plus swap: about $17\times$
short, combined, down from a standing bracket of $2.4$--$195$~TB of storage.
The gap has narrowed by roughly three orders of magnitude; it has not closed.
Every level-7 figure here is an extrapolation --- the run was not performed ---
and the level-6-to-7 state multiplier is the one input for which no cheaper
rehearsal exists, precisely because $\nmin(7) = 13$.

\paragraph{Why splitting the problem does not help.}
The obvious response is to decompose: partition the search by comparator
prefix, run $k$ independent jobs, and pay $k \cdot (N/k)^{e}$ instead of
$N^{e}$. Across $146$ runs --- $144$ prefix jobs at six prefix depths and two
levels, plus two monolithic controls --- the lowest peak per-job memory ever
observed was $0.990\times$ the monolithic run's, and that was at depth~1, where
$k=1$ and the split is a no-op. The best non-trivial point is $0.998\times$; by
depth~4 the peak per-job memory is $23.6\times$ the monolith's. Level~7 would
need $0.095\times$.

The failure is at the first premise, not in the arithmetic: $k \cdot (N/k)^e$
assumes the split partitions the state space, and prefix cones \emph{cover} it.
At level~5, depth~2, one job's memo held $11{,}849{,}694$ entries against the
entire monolithic run's $11{,}629{,}126$ --- $1.9\%$ \emph{larger} than the run
it was supposed to be one third of. A memo diff between two sibling cones puts
the job-private fraction at $36.2\%$ of entries, so $63.8\%$ of the work is
shared and is paid $k$ times over. The mechanism is general: an
iterative-deepening branch-and-bound never solves an interior subproblem to
completion --- it raises the root's bound incrementally, improving each
successor only far enough to move the current minimum --- whereas each
decomposed job must certify its own root to its full value. The parts are
strictly harder than the roles those same parts play inside the whole. Granting
the pro-decomposition argument in full --- every wide state perfectly private,
the low-width remainder free --- the census supplies a floor: widths $\ge 10$
are $56.7\%$ of packed bytes, so even a perfect decomposition with
$k \to \infty$ and zero CPU cost leaves the memo at $0.43$ times its monolithic
size, still several times over the available memory. This is a clean negative
and it ends that line of work.

%======================================================================
\section{Trusted base}\label{sec:trust}
%======================================================================

Following Harder~\cite{Harder20}, we do not claim the components below are free
of bugs and do not expect them to be. The goal is narrower: to say which could,
if wrong, move a stated claim. The scope is the specific runs reported here, so
only bugs triggered during those runs are relevant, and we adopt the de Bruijn
criterion~\cite{BW05} as the standard the certified computation is meant to
meet rather than claiming correctness directly.

\begin{center}
\begin{tabular}{>{\raggedright\arraybackslash}p{0.28\textwidth}c>{\raggedright\arraybackslash}p{0.52\textwidth}}
\toprule
component & trusted? & note \\
\midrule
Isabelle/HOL kernel (Isabelle2020) & Yes &
  Shared with the checker development, so the two formal artifacts are not
  independent of each other at this level. \\
Statements in \texttt{Ambient\_Collapse.thy} & Yes &
  The proofs are machine-checked; a mis-\emph{stated} theorem is not. Two of
  our hypotheses were mis-stated and the mechanization caught them
  (Section~\ref{sec:corrections}); a third would not be. \\
Proofs in \texttt{Ambient\_Collapse.thy} & No &
  Machine-checked, no \texttt{sorry}/\texttt{oops}/\texttt{axiomatization}. \\
That the sorting-network search satisfies H0--H5 & Yes &
  \textbf{Not machine-checked.} Rests on an exhaustive reading of the
  recursion plus a second reviewer given the code as a refutation target; the
  one leak found is Section~\ref{sec:instance}. This is the largest single
  trusted item in the paper. \\
Our rebuilt search engine & No &
  Its output is a certificate, checked independently. \\
Prune-all and certificate assembly & No &
  Same; an error here yields a rejected certificate. \\
Extracted verified checker core & Yes &
  An \emph{unchanged} extraction of the frozen \texttt{Checker.thy}; the
  extraction pipeline was reproduced from source under Isabelle2020. \\
Checker wrapper and decoder & Yes &
  \textbf{Patched}, and not verified: a large-read fix in the Haskell wrapper
  was required on the verification host. ``Unmodified'' scopes to the extracted
  core, never to the binary. \\
Census figures at levels 5 and 6 & Yes &
  One measured ambient each, so the affine law there rests on width-stratum
  evidence rather than an observed common difference; $|\Reach(13,6)|$ is a
  prediction of Theorem~\ref{thm:partII}. \\
Certificate step and memory figures & No &
  No claim depends on them; the step count is a recount over the certificate,
  superseding an earlier figure derived from a different quantity. \\
Hypothesis checks in the two further domains & Yes &
  Python, not formalized. In the Boolean domain a raw enumeration independently
  checks the canonical form; in prefix reversal there is no such check, and H4
  is verified over a finite range, not proved. \\
GHC, operating system, hardware & Yes & Standard, and unavoidable. \\
\bottomrule
\end{tabular}
\end{center}

Two dependencies deserve naming before a referee names them. The mechanized
theorem and the certified computation are naturally described as two
independent formal artifacts; they share the Isabelle kernel and the same
released version of it, so they are independent in content but not in
toolchain. And that toolchain is the programme's most fragile asset: the
Isabelle2020 installation lives outside the repository, has had to be
re-downloaded once already, and the theory sources are committed while the
prover that checks them is not pinned. Of the four stages of
Section~\ref{sec:s11}, only verification left a machine-readable artifact.

Finally, a deliberate exclusion. Our own abstract theory sits \emph{outside}
the trusted base for the certified bound: every line of
Sections~\ref{sec:class}--\ref{sec:mech} could be wrong and $S(11) \ge 35$
would not move, because that bound is established by a certificate and a
checker that know nothing about ambients, levels or censuses. Conversely, the
collapse theorem is what licenses the claim that one ambient-11 census
determines the ambient-13 one, and \emph{that} rests on the audit in the fourth
row above --- the first place a reader who wants to disbelieve us should look.

%======================================================================
\section{Conclusion}
%======================================================================

An exhaustive combinatorial-optimality search can have a cutoff in the sense
that parameterized verification has understood the word for thirty years: above
a computable threshold, the object being computed at parameter $n$ is
determined by the object at the threshold, plus a closed-form correction. We
have given hypotheses under which this holds for the census --- the set of
states such a search stores --- proved both halves and the growth law in
Isabelle/HOL with two interpretations that make them non-vacuous, and
instantiated the class in three domains, one of which is a positive control and
one of which was chosen because it could have falsified the theorem and did
not. The growth law we prove is affine with slope exactly one; the polynomial
version is a machine-checked implication whose hypothesis is measured rather
than derived. What we have not shown is that any particular engine is an
instance: that remains an audit, and it is the obligation a reader should press
on. The method's practical reach ends where the theorem stops supplying a
smaller ambient to collapse to --- and there, as we measured, the cost is
governed by a quantity the theorem does not control.

\footnotesize
\bibliographystyle{plain}
\bibliography{refs}

\end{document}
